\documentclass[11pt]{article}
\usepackage[T1]{fontenc}
\usepackage[utf8]{inputenc}
\usepackage[french]{babel}
\usepackage{lmodern}
\usepackage{geometry}
\geometry{margin=1in}
\usepackage{microtype}
\usepackage{amsmath,amssymb,mathtools,mathrsfs}
\usepackage{enumitem}
\usepackage{amsthm}
\usepackage{tikz-cd}
\DeclareMathOperator{\Reg}{Reg}
\DeclareMathOperator{\Supp}{Supp}
\DeclareMathOperator{\Div}{Div}
\DeclareMathOperator{\codim}{codim}
\providecommand{\cO}{\mathcal{O}}
\providecommand{\fm}{\mathfrak{m}}
\providecommand{\wt}[1]{\widetilde{#1}}
\providecommand{\Cons}{\operatorname{Cons}}
\providecommand{\Hh}{\mathrm{H}}
\providecommand{\uH}{\underline{\mathrm H}}
\providecommand{\uExt}{\underline{\Ext}}
\providecommand{\qedtext}{\hfill$\square$}
\providecommand{\ZZ}{\mathbb{Z}}
\providecommand{\QQ}{\mathbb{Q}}
\providecommand{\RR}{\mathbb{R}}
\providecommand{\FF}{\mathbb{F}}
\providecommand{\NN}{\mathbb{N}}

\newenvironment{statement}[1]{\par\medskip\noindent\textbf{#1.}\ }{\par\medskip}
\newenvironment{prooftext}[1][Proof]{\par\noindent\emph{#1.}\ }{\par\medskip}

\usepackage{hyperref}
\hypersetup{colorlinks=true,linkcolor=blue,citecolor=blue,urlcolor=blue}
\setlength{\parindent}{0pt}
\setlength{\parskip}{0.55em}
\let\accentH\H
\renewcommand{\H}{\mathrm{H}}
\DeclareMathOperator{\Hom}{Hom}
\DeclareMathOperator{\Ext}{Ext}
\DeclareMathOperator{\Tor}{Tor}
\DeclareMathOperator{\Spec}{Spec}
\DeclareMathOperator{\cd}{cd}
\newcommand{\R}{\mathrm{R}}
\newcommand{\D}{\mathrm{D}}
\newcommand{\Z}{\mathbb{Z}}
\newcommand{\F}{\mathbb{F}}
\newcommand{\Q}{\mathbb{Q}}
\newcommand{\E}{\mathrm{E}}
\newcommand{\ob}{\operatorname{ob}}
\newcommand{\car}{\operatorname{car}}
\newcommand{\ol}[1]{\overline{#1}}
\newcommand{\SGAA}{SGA~A}

\providecommand{\Ker}{\operatorname{Ker}}
\providecommand{\Coker}{\operatorname{Coker}}
\providecommand{\Gal}{\operatorname{Gal}}
\providecommand{\car}{\operatorname{car}}


\providecommand{\RHom}{\mathrm{RHom}}
\providecommand{\Dctf}{D_{\mathrm{ctf}}}
\providecommand{\Id}{\mathrm{Id}}
\providecommand{\pr}{\operatorname{pr}}
\providecommand{\Gr}{\Gamma}
\providecommand{\Fix}{\operatorname{Fix}}
\providecommand{\Loc}{\operatorname{Loc}}
\providecommand{\lto}{\longrightarrow}
\providecommand{\xto}[1]{\xrightarrow{#1}}
\providecommand{\lisom}{\xrightarrow{\sim}}
\providecommand{\ctf}{\mathrm{ctf}}
\providecommand{\et}{\mathrm{\acute{e}t}}
\providecommand{\Sfrac}{SGA~$4^{1/2}$}


\providecommand{\cl}{\operatorname{cl}}
\providecommand{\Tr}{\operatorname{Tr}}
\providecommand{\EP}{\operatorname{EP}}
\providecommand{\gr}{\operatorname{gr}}
\providecommand{\DF}{\operatorname{DF}}

\providecommand{\OO}{\mathcal{O}}
\providecommand{\LL}{\mathbf L}
\providecommand{\parf}{\mathrm{parf}}
\providecommand{\coh}{\mathrm{coh}}
\providecommand{\frob}{\operatorname{Frob}}
\nofiles

\providecommand{\RGamma}{\mathrm R\Gamma}
\providecommand{\Rj}{\mathrm Rj}
\providecommand{\Tr}{\operatorname{Tr}}
\providecommand{\id}{\operatorname{id}}
\providecommand{\cone}{\operatorname{cone}}
\providecommand{\lten}{\overset{\mathrm L}{\otimes}}
\providecommand{\Dctf}{D_{\mathrm{ctf}}}
\providecommand{\Ob}{\operatorname{ob}}


\providecommand{\Tr}{\operatorname{Tr}}
\providecommand{\RHom}{\mathrm{RHom}}
\providecommand{\Hom}{\operatorname{Hom}}
\providecommand{\ob}{\operatorname{ob}}
\providecommand{\lten}{\overset{\mathrm L}{\otimes}}
\providecommand{\nat}{\natural}
\providecommand{\dfn}{\mathrm{dfn}}


\providecommand{\Fl}{\operatorname{Fl}}
\providecommand{\MLAR}{\mathrm{MLAR}}
\providecommand{\AR}{\mathrm{AR}}
\providecommand{\NN}{\mathbb N}
\providecommand{\Ens}{\mathrm{Ens}}
\providecommand{\cC}{\mathcal C}
\providecommand{\cE}{\mathcal E}
\providecommand{\cF}{\mathcal F}
\providecommand{\cG}{\mathfrak G}
\providecommand{\source}{\operatorname{source}}
\providecommand{\op}{\mathrm{op}}


\providecommand{\End}{\operatorname{End}}
\providecommand{\id}{\operatorname{id}}
\providecommand{\Im}{\operatorname{Im}}
\providecommand{\Ker}{\operatorname{Ker}}
\providecommand{\Coker}{\operatorname{Coker}}
\providecommand{\gr}{\operatorname{gr}}

\renewcommand{\Im}{\operatorname{Im}}

\providecommand{\Rlim}{\varprojlim}
\providecommand{\adn}{\operatorname{adn}}
\providecommand{\adf}{\operatorname{adf}}
\providecommand{\AR}{\mathrm{AR}}
\providecommand{\MLAR}{\mathrm{MLAR}}
\providecommand{\Ecat}{\mathcal E}
\providecommand{\Dcat}{\mathcal D}
\providecommand{\Ccat}{\mathcal C}
\providecommand{\Bcat}{\mathcal B}
\providecommand{\Pcat}{\mathcal P}
\providecommand{\J}{J}
\providecommand{\gr}{\operatorname{gr}}
\providecommand{\grs}{\operatorname{grs}}
\providecommand{\Homcat}{\operatorname{Hom}}
\providecommand{\Ker}{\operatorname{Ker}}
\providecommand{\Coker}{\operatorname{Coker}}
\providecommand{\Im}{\operatorname{Im}}
\providecommand{\can}{\operatorname{can}}
\providecommand{\epi}{\operatorname{epi}}
\providecommand{\id}{\operatorname{id}}


\providecommand{\RGamma}{\mathrm R\Gamma}
\providecommand{\RHom}{\mathrm R\mathcal{H}om}
\providecommand{\Zell}{\mathbb Z_\ell}
\providecommand{\Qell}{\mathbb Q_\ell}
\providecommand{\Ql}{\mathbb Q_\ell}
\providecommand{\lten}{\overset{\mathrm L}{\otimes}}
\providecommand{\Ob}{\operatorname{ob}}
\providecommand{\Ens}{\operatorname{Ens}}
\providecommand{\Ensf}{\operatorname{Ensf}}
\providecommand{\Aut}{\operatorname{Aut}}
\providecommand{\End}{\operatorname{End}}
\providecommand{\Sp}{\operatorname{Sp}}
\providecommand{\Ab}{\operatorname{Ab}}
\providecommand{\Abc}{\operatorname{Abc}}
\providecommand{\Abf}{\operatorname{Abf}}
\providecommand{\fct}{\operatorname{fct}}
\providecommand{\Lc}{\operatorname{Lc}}
\providecommand{\GL}{\operatorname{GL}}
\providecommand{\car}{\operatorname{car}}
\providecommand{\cHom}{\mathcal{H}om}
\providecommand{\fc}{\operatorname{fc}}
\providecommand{\adc}{\operatorname{adc}}
\providecommand{\ad}{\operatorname{ad}}
\providecommand{\adn}{\operatorname{adn}}
\providecommand{\modn}{\operatorname{mod}}
\providecommand{\AR}{\operatorname{AR}}


\providecommand{\R}{\mathrm R}
\providecommand{\V}{\mathrm V}
\providecommand{\A}{\mathbb A}
\providecommand{\P}{\mathbb P}
\providecommand{\Z}{\mathbb Z}
\providecommand{\muS}{\mu_S}
\providecommand{\cupdot}{\smile}
\providecommand{\id}{\operatorname{id}}
\providecommand{\Tr}{\operatorname{Tr}}
\providecommand{\Spec}{\operatorname{Spec}}
\providecommand{\Supp}{\operatorname{Supp}}
\providecommand{\coker}{\operatorname{coker}}
\providecommand{\lten}{\overset{\mathrm L}{\otimes}}
\newenvironment{proposition}[1][]{\par\medskip\noindent\textbf{Proposition #1.}\ }{\par\medskip}
\newenvironment{corollaire}[1][]{\par\medskip\noindent\textbf{Corollaire #1.}\ }{\par\medskip}
\newenvironment{theoreme}[1][]{\par\medskip\noindent\textbf{Théorème #1.}\ }{\par\medskip}
\newenvironment{lemme}[1][]{\par\medskip\noindent\textbf{Lemme #1.}\ }{\par\medskip}


\providecommand{\Ahat}{\widehat A}
\providecommand{\Gm}{\mathbb G_m}
\providecommand{\Et}{\mathrm{et}}
\providecommand{\Pic}{\operatorname{Pic}}
\providecommand{\cl}{\operatorname{cl}}
\providecommand{\EP}{\operatorname{EP}}
\providecommand{\Ch}{\operatorname{Ch}}
\providecommand{\lten}{\overset{\mathrm L}{\otimes}}
\providecommand{\rank}{\operatorname{rank}}
\providecommand{\Sym}{\operatorname{Sym}}
\providecommand{\Spec}{\operatorname{Spec}}
\providecommand{\pr}{\operatorname{pr}}
\providecommand{\id}{\operatorname{id}}
\providecommand{\im}{\operatorname{Im}}
\providecommand{\coker}{\operatorname{Coker}}


\providecommand{\cC}{\mathcal C}
\providecommand{\cD}{\mathcal D}
\providecommand{\cA}{\mathcal A}
\providecommand{\cB}{\mathcal B}
\providecommand{\Ab}{\operatorname{Ab}}
\providecommand{\Ob}{\operatorname{Ob}}
\providecommand{\parf}{\mathrm{parf}}
\providecommand{\coh}{\mathrm{coh}}
\providecommand{\lten}{\overset{\mathrm L}{\otimes}}


\providecommand{\Mod}{\operatorname{Mod}}
\providecommand{\tr}{\operatorname{tr}}
\providecommand{\Sw}{\operatorname{Sw}}
\providecommand{\Art}{\operatorname{Art}}
\providecommand{\etG}{\mathrm{\acute et},G}
\providecommand{\cT}{\mathcal T}


\providecommand{\clK}[1]{\operatorname{cl}_{K(#1)}}
\providecommand{\clKstar}[1]{\operatorname{cl}_{K^*(#1)}}
\providecommand{\epsD}{\varepsilon^\Delta}
\providecommand{\alphaD}{\alpha^\Delta}
\providecommand{\bbone}{\mathbf 1}
\providecommand{\sw}{\operatorname{sw}}
\providecommand{\art}{\operatorname{art}}
\providecommand{\RG}{\mathrm R\Gamma}
\providecommand{\RGc}{\mathrm R\Gamma_c}
\providecommand{\RHom}{\mathrm{RHom}}
\providecommand{\Rpi}{\mathrm R\pi}
\providecommand{\rG}{\mathrm R\Gamma^G}
\providecommand{\wh}[1]{\widehat{#1}}
\providecommand{\checkSw}{\check{\operatorname{Sw}}}


\providecommand{\La}{\Lambda}
\providecommand{\bZ}{\mathbb Z}
\providecommand{\bQ}{\mathbb Q}
\providecommand{\bZell}{\mathbb Z_\ell}
\providecommand{\fm}{\mathfrak m}
\providecommand{\card}{\operatorname{card}}
\providecommand{\im}{\operatorname{im}}
\providecommand{\Aut}{\operatorname{Aut}}
\providecommand{\Hom}{\operatorname{Hom}}
\providecommand{\Spec}{\operatorname{Spec}}
\providecommand{\cl}{\operatorname{cl}}
\providecommand{\NW}{\operatorname{NW}}
\providecommand{\Tr}{\operatorname{Tr}}
\providecommand{\Ob}{\operatorname{ob}}
\providecommand{\car}{\operatorname{car}}
\providecommand{\qedtext}{\hfill$\square$}


\providecommand{\fr}{\operatorname{fr}}
\providecommand{\Fr}{\operatorname{Fr}}
\providecommand{\Sch}{\operatorname{Sch}}
\providecommand{\Proj}{\operatorname{Proj}}
\providecommand{\Sym}{\operatorname{Sym}}
\providecommand{\PP}{\mathbb P}
\providecommand{\JJ}{\mathcal J}
\providecommand{\XX}{\mathcal X}
\providecommand{\cL}{\mathcal L}
\providecommand{\cM}{\mathcal M}
\providecommand{\cN}{\mathcal N}
\providecommand{\cS}{\mathcal S}
\providecommand{\cX}{\mathcal X}
\providecommand{\cY}{\mathcal Y}
\providecommand{\cZ}{\mathcal Z}
\providecommand{\bF}{\mathbb F}
\providecommand{\bbZ}{\mathbb Z}
\providecommand{\simeqto}{\xrightarrow{\sim}}



\providecommand{\Fr}{\operatorname{Fr}}
\providecommand{\fr}{\operatorname{fr}}
\providecommand{\Sch}{\operatorname{Sch}}
\providecommand{\Spec}{\operatorname{Spec}}
\providecommand{\id}{\operatorname{id}}
\providecommand{\RGamma}{\mathrm R\Gamma}
\providecommand{\Ob}{\operatorname{Ob}}
\providecommand{\Gal}{\operatorname{Gal}}
\providecommand{\et}{\acute{e}t}
\providecommand{\lten}{\overset{\mathrm L}{\otimes}}
\providecommand{\boxt}{\boxtimes}
\begin{document}

\subsection*{4. Morphisme de Frobenius itéré; préschémas parfaits}

Pour tout entier $\nu\geq 0$, l'endomorphisme itéré $\fr_X^\nu$ dépend fonctoriellement du préschéma $X$ de caractéristique $p$. Si donc $S$ est un préschéma de caractéristique $p$ et $X$ un $S$-préschéma, $\fr_X^\nu$ se factorise, d'une manière unique, en un $S$-morphisme
\[
\Fr_{X/S}^\nu:X\longrightarrow X^{(p^\nu/S)},
\]
où
\[
X^{(p^\nu/S)}=X\times_S(S,\fr_S^\nu),
\]
suivi de la projection $X^{(p^\nu/S)}\to X$. Le morphisme $\Fr_{X/S}^\nu$ dépend fonctoriellement du $S$-préschéma $X$ et il est compatible avec le changement de la base $S$ (cf. n$^\circ$2, prop. 1b)). D'ailleurs il est clair que
\[
X^{(p^\nu/S)}=X^{(p^\nu)}\simeq\bigl(X^{(p^{\nu-1})}\bigr)^{(p)}
\]
et que
\[
\Fr_{X/S}^\nu=(\Fr_{X/S}^{\nu-1})^{(p)}\circ\Fr_{X/S},
\]
modulo l'identification permise par l'isomorphisme canonique précédent.

\begin{statement}{Définition 4}
On dit qu'un préschéma $S$ de caractéristique $p$ est parfait si $\fr_S$ est un automorphisme de $S$.
\end{statement}

Si $S$ est un préschéma de caractéristique $p$ parfait et $X$ un $S$-préschéma, on voit immédiatement que $X^{(p/S)}$ s'identifie au préschéma $X$ considéré comme $S$-préschéma au moyen du morphisme $\fr_S^{-1}\circ g:X\to S$, où $g$ est le morphisme structural de $X$ dans $S$; modulo cette identification, $\Fr_{X/S}=\fr_X$.

Comme exemple de préschéma parfait, nous rencontrerons $S=\Spec(\mathbb F_q)$, où $q=p^\nu$ est une puissance de la caractéristique $p$; en effet, dans ce cas $\fr_S^\nu=\id_S$. Si $X$ est un $S$-préschéma, on voit que
\[
X^{(p^\nu/S)}=X,
\qquad
\Fr_{X/S}^\nu=\fr_X^\nu.
\]
Notons que pour $\nu=1$ on obtient $S=\Spec(\mathbb F_p)$, pour lequel $\fr_S=\id_S$; les $S$-préschémas sont les préschémas de caractéristique $p$; sur cette base $S$, $X^{(p/S)}=X$ et $\Fr_{X/S}=\fr_X$.

\section*{§ 2. Correspondance de Frobenius}

\subsection*{1. Définition}

Soient $X$ un préschéma de caractéristique $p$ et $F$ un faisceau d'ensembles sur $X_{\et}$. Pour tout $X$-préschéma étale $U$ on a
\[
(\fr_X)_*(F)(U)=F(\fr_X^{-1}(U))=F(U^{(p/X)}),
\]
et le morphisme de Frobenius $\Fr_{U/X}:U\to U^{(p/X)}$ définit une application
\[
F(\Fr_{U/X}):\fr_{X*}(F)(U)\longrightarrow F(U),
\]
visiblement fonctorielle en $U$ (puisque $\Fr_{U/X}$ est fonctoriel en $U$). Ainsi on obtient un morphisme de faisceaux sur $X_{\et}$:
\[
(\fr_X)_*(F)\longrightarrow F.
\]
Or, lorsque $U$ est étale sur $X$, $\Fr_{U/X}$ est un isomorphisme (§ 1 n$^\circ$2, prop. 2c)), donc il en est de même de $F(\Fr_{U/X})$; il en résulte que le morphisme précédent est un isomorphisme. Dans le cas où $F$ est muni d'une structure algébrique (par exemple si c'est un faisceau de $\Lambda$-modules, $\Lambda$ anneau de base donné), cet isomorphisme est compatible avec la structure algébrique. L'inverse
\[
F(\Fr_{/X})^{-1}:F\longrightarrow \fr_{X*}(F)
\]
de cet isomorphisme définit, puisque $(\fr_X)^*$ est adjoint à gauche du foncteur $(\fr_X)_*$, un morphisme
\[
\Fr^*_{F/X}:(\fr_X)^*(F)\longrightarrow F.
\]
On peut même dire plus, car $\fr_X$ étant un homéomorphisme universel (§ 1 n$^\circ$2, prop. 2a)), les foncteurs $(\fr_X)^*$ et $(\fr_X)_*$ sont quasi-inverses l'un de l'autre, et par suite $\Fr^*_{F/X}$ est un isomorphisme comme le morphisme $F(\Fr_{/X})^{-1}$ dont il provient.

\begin{statement}{Définition 1}
Étant donnés un préschéma $X$ de caractéristique $p$ et un faisceau d'ensembles $F$ sur $X_{\et}$, on appelle correspondance de Frobenius sur $(X,F)$ le couple $(\fr_X,\Fr^*_{F/X})$ formé de l'endomorphisme de Frobenius $\fr_X$ de $X$ et de l'isomorphisme
\[
\Fr^*_{F/X}=(F(\Fr_{/X})^{-1})^\sharp:(\fr_X)^*(F)\longrightarrow F.
\]
\end{statement}

Bien entendu on définit de la même façon une classe de correspondance de Frobenius itérée $(\fr_X^\nu,(\Fr^*_{F/X})^\nu)$ pour tout entier $\nu\geq0$; l'isomorphisme
\[
(\Fr^*_{F/X})^\nu:(\fr_X^\nu)^*(F)\longrightarrow F
\]
est égal à
\[
(\Fr^*_{F/X})^{\nu-1}\circ (\fr_X^{\nu-1})^*(\Fr^*_{F/X})
\]
ou encore à
\[
\Fr^*_{F/X}\circ\fr_X^*((\Fr^*_{F/X})^{\nu-1})
\]
pour $\nu\geq1$.

Examinons le cas particulier où $F$ est représentable par un préschéma $V$ étale sur $X$; alors $(\fr_X)^*(F)=V^{(p/X)}$, $F(V)=\Hom_X(V,V)$,
\[
(\fr_X)^*(F)(V)=F(V^{(p/X)})=\Hom_X(V^{(p/X)},V),
\]
et $F(\Fr_{V/X})$ est l'application de $\Hom_X(V^{(p/X)},V)$ dans $\Hom_X(V,V)$ qui transforme un morphisme $\psi$ en $\psi\circ\Fr_{V/X}$; or $\Fr^*_{F/X}$ s'identifie à un morphisme de $V^{(p/X)}$ dans $V$ qui est l'image de $\id_V$ par l'application $F(\Fr_{V/X})^{-1}$ réciproque de la précédente. Ceci montre que pour $F=V$ étale sur $X$ on a
\[
\Fr^*_{V/X}=(\Fr_{V/X})^{-1}.
\]

Notons que si le faisceau $F$ est constant il est clair que
\[
(\fr_X)^*(F)=F,
\qquad
\Fr^*_{F/X}=\id_F.
\]

\begin{statement}{Proposition 1}
\begin{enumerate}[label=\alph*)]
\item L'isomorphisme $\Fr^*_{F/X}$ dépend fonctoriellement du faisceau $F$; autrement dit on a un isomorphisme fonctoriel
\[
\Fr^*_{/X}\,\fr_X^*\longrightarrow \id_{\widehat X_{\et}},
\]
où $\widehat X_{\et}$ est la catégorie des faisceaux d'ensembles sur $X_{\et}$.
\item La correspondance de Frobenius $(\fr_X,\Fr^*_{F/X})$ sur $(X,F)$ est compatible avec le changement de la base $X$. C'est-à-dire que, si $\mathcal F$ désigne la catégorie fibrée sur $(\Sch)_p$ des faisceaux, pour la topologie étale, sur des préschémas de caractéristique $p$ variables, les foncteurs $\fr_X^*$, $X\in\Ob(\Sch)_p$, définissent un $(\Sch)_p$-foncteur cartésien
\[
\fr^*: \mathcal F\longrightarrow\mathcal F,
\]
et les morphismes $\Fr^*_{F/X}$, $X\in\Ob(\Sch)_p$, $F\in\Ob\widehat X_{\et}$, définissent un $(\Sch)_p$-morphisme fonctoriel
\[
\Fr^*:\fr^*\longrightarrow\id_{\mathcal F}
\]
qui est un isomorphisme.
\item Inversement ces propriétés caractérisent $\Fr^*$, c'est-à-dire qu'il existe un $(\Sch)_p$-morphisme fonctoriel et un seul de $\fr^*$ dans $\id_{\mathcal F}$.
\end{enumerate}
\end{statement}

Remarquons que l'énoncé b) implique l'énoncé a); ce dernier est tout à fait évident.

Démontrons b). Soient $g:X'\to X$ un morphisme de préschémas de caractéristique $p$ et $F$ un faisceau sur $X_{\et}$; posons $F'=g^*(F)$ et, plus généralement, affectons d'un prime l'image inverse par $g$ de tout objet donné sur $X$. Comme
\[
g\circ\fr_{X'}=\fr_X\circ g,
\]
on a
\[
\fr_{X'}^*(F')=\fr_{X'}^*(g^*F)=g^*(\fr_X^*F)=(\fr_X^*F)'
\]
d'où le fait que les $\fr_X^*$ définissent un foncteur cartésien $\fr^*$ au-dessus de $(\Sch)_p$. Il reste à prouver que
\[
\Fr^*_{F'/X'}=(\Fr^*_{F/X})'.
\]
En utilisant les propriétés des foncteurs adjoints, on voit que cela revient à montrer que l'image de $F(\Fr_{/X'})^{-1}$ par l'application bijective
\[
\Hom_{X'}(F',\fr_{X'*}(F'))
 =\Hom_{X'}(g^*F,\fr_{X'*}g^*F)
 \simeq \Hom_X(F,g_*\fr_{X'*}g^*F)
\]
coïncide avec l'image de $F(\Fr_{/X})^{-1}$ par l'application
\[
\Hom_X(F,\fr_{X*}F)\longrightarrow
\Hom_X(F,g_*\fr_{X'*}g^*F)
\]
définie par le morphisme d'adjonction $F\to g_*g^*F$; autrement dit, il faut montrer que
\[
\alpha\circ F(\Fr_{/X})=g_*(F'(\Fr_{/X'}))\circ\fr_{X*}(\alpha),
\]
où $\alpha:F\to g_*g^*F$ est l'adjonction. C'est-à-dire que, pour tout préschéma $U$ étale sur $X$, on a un diagramme commutatif
\[
\begin{tikzcd}[column sep=large,row sep=large]
F(U^{(p/X)}) \arrow[r,"F(\Fr_{U/X})"] \arrow[d]
& F(U) \arrow[d]\\
F'(U'{}^{(p/X')}) \arrow[r,"F'(\Fr_{U'/X'})"']
& F'(U') ,
\end{tikzcd}
\]
où les flèches verticales sont les applications canoniques; or ceci provient immédiatement de
\[
\Fr_{U'/X'}=(\Fr_{U/X})'
\]
(§ 1, n$^\circ$2, prop. 1b)).

Pour prouver c), on se ramène, en composant avec $(\Fr^*)^{-1}$, à prouver que $\id_{\mathcal F}$ a un seul $(\Sch)_p$-endomorphisme: l'identité. Or, si $\psi$ est un $(\Sch)_p$-endomorphisme de $\id_{\mathcal F}$, soient $F$ un faisceau sur un préschéma $X$ de caractéristique $p$ et $U$ un préschéma étale sur $X$; on sait que $F(U)$ s'identifie à $\Hom_U(U,F|U)$ et $\psi_F(U):F(U)\to F(U)$ à l'application
\[
s\longmapsto \psi_{F|U}\circ s=s\circ\psi_U,
\]
où $s\in\Hom_U(U,F|U)$; comme le faisceau final $U$ de $\widetilde U_{\et}$ n'admet pas d'autre endomorphisme que l'identité, on a $\psi_U=\id_U$, il en résulte que $\psi_F(U)$ est aussi l'identité de $F(U)$ et que $\psi=\id_{\mathcal F}$.

Bien entendu, si le faisceau $F$ est muni d'une structure algébrique, l'isomorphisme $\Fr^*_{F/X}$ est compatible avec cette structure; par exemple si $F$ est un faisceau de groupes sur le préschéma $X$ de caractéristique $p$, $\Fr^*_{F/X}$ est un isomorphisme de faisceaux de groupes. Soit $\Lambda$ un anneau et soit $F$ un faisceau de $\Lambda$-modules sur $X$; alors $\Fr^*_{F/X}$ est un isomorphisme de faisceaux de $\Lambda$-modules. Plus généralement, si $F^\bullet$ est un complexe de $\Lambda$-modules, on définit immédiatement un isomorphisme
\[
\Fr^*_{F^\bullet/X}:\fr_X^*(F^\bullet)\longrightarrow F^\bullet
\]
égal à $\Fr^*_{F^i/X}$ sur la composante de degré $i$, grâce à la proposition 1a). Dans ce cas le couple $(\fr_X,\Fr^*_{F^\bullet/X})$ est bien une correspondance au sens habituel sur $(X,F^\bullet)$; cette correspondance est fonctorielle en $F^\bullet$ et compatible avec le changement de la base $X$.

\subsection*{2. Effet sur la cohomologie}

\begin{statement}{Proposition 2}
Soit $X$ un préschéma de caractéristique $p$.
\begin{enumerate}[label=\alph*)]
\item Pour tout faisceau d'ensembles $F$ sur $X_{\et}$, l'endomorphisme de $H^0(X,F)$ défini par $(\fr_X,\Fr^*_{F/X})$, c'est-à-dire le composé
\[
H^0(X,F)\longrightarrow H^0(X,\fr_X^*F)\longrightarrow H^0(X,F),
\]
où la première flèche est l'application canonique, caractère contravariant de $H^0(X,F)$ par rapport à $X$, et la seconde est $H^0(X,\Fr^*_{F/X})$, est l'identité de $H^0(X,F)$.
\item Pour tout faisceau de groupes $F$ sur $X_{\et}$, l'endomorphisme de $H^1(X,F)$ défini par $(\fr_X,\Fr^*_{F/X})$, c'est-à-dire
\[
H^1(X,F)\longrightarrow H^1(X,\fr_X^*F)\longrightarrow H^1(X,F),
\]
est l'identité.
\item Pour tout complexe limité à gauche $F$ de $\Lambda$-modules sur $X_{\et}$, l'endomorphisme de $\RGamma(X,F)$ défini par $(\fr_X,\Fr^*_{F/X})$, c'est-à-dire le composé
\[
\RGamma(X,F)\longrightarrow \RGamma(X,\fr_X^*F)\longrightarrow \RGamma(X,F),
\]
est l'identité.
\end{enumerate}
\end{statement}

Établissons d'abord a); nous en déduirons b) et c). Soit $\psi_F$ l'endomorphisme considéré de $H^0(X,F)$; il est clair que $\psi_F$ dépend fonctoriellement de $F$ et que $\psi_X=\id_{H^0(X,X)}$. Or tout élément $s$ de $H^0(X,F)$ s'identifie à un morphisme de faisceaux $s:X\to F$ tel que $s$ soit l'image par $H^0(X,s)$ de l'unique élément de $H^0(X,X)$; comme
\[
\psi_F\circ H^0(X,s)=H^0(X,s)\circ\psi_X=H^0(X,s),
\]
d'après les remarques précédentes, on voit que $\psi_F(s)=s$ pour tout $s$, donc $\psi_F=\id_{H^0(X,F)}$.

Notons que $\psi_F$ est encore le composé
\[
H^0(X,F)\longrightarrow H^0(X,\fr_{X*}F)
\longrightarrow H^0(X,F),
\]
où la première flèche est $H^0(X,F(\Fr_{/X})^{-1})$ et la seconde l'application canonique, qui n'est autre que l'identité. De même les endomorphismes considérés en b) et c) se factorisent ainsi:
\[
H^1(X,F)\longrightarrow H^1(X,\fr_{X*}F)
\longrightarrow H^1(X,F),
\]
et
\[
\RGamma(X,F)\longrightarrow \RGamma(X,\fr_{X*}F)
\longrightarrow \RGamma(X,F).
\]

Pour établir b), nous utiliserons un monomorphisme $F\to G$ de $F$ dans un faisceau de groupes $G$ tel que $H^1(X,G)=0$; alors $\fr_{X*}F\to\fr_{X*}G$ est un monomorphisme et $H^1(X,\fr_{X*}G)=0$. Introduisons le faisceau d'ensembles homogènes quotient $C=G/F$. Comme $\fr_X$ est entier, surjectif et radiciel, les foncteurs $\fr_X^*$ et $\fr_{X*}$ sont des équivalences de catégories quasi-inverses l'une de l'autre (SGA VIII théorème 1.1); on en déduit que
\[
\fr_{X*}(C)\simeq\fr_{X*}(G)/\fr_{X*}(F),
\]
et la suite exacte de cohomologie (SGA XII proposition 3.1) montre que l'endomorphisme de $H^1(X,F)$ considéré provient par passage au quotient de
\[
\psi_C:H^0(X,C)\longrightarrow H^0(X,C),
\]
qui est l'identité d'après a); cet endomorphisme est donc $\id_{H^1(X,F)}$.

L'assertion c) se démontre d'une manière analogue, en utilisant un complexe $I$ injectif en chaque degré et isomorphe à $F$ dans la catégorie dérivée $D^+$, construite sur les $\Lambda$-modules; le complexe $\fr_{X*}(I)$ est injectif en chaque degré et isomorphe à $\fr_{X*}(F)$ dans $D^+$, car $\fr_{X*}$ est une équivalence de catégories. On peut alors identifier $\RGamma(X,F)$ à $H^0(X,I)$ et $\RGamma(X,\fr_X^*F)$ à $H^0(X,\fr_{X*}I)$; on voit que l'endomorphisme considéré de $\RGamma(X,F)$ s'identifie à $\psi_I$, qui est l'identité; cet endomorphisme est donc $\id_{\RGamma(X,F)}$.

\subsection*{3. Comportement des produits}

Soient $X$ et $X'$ des préschémas de caractéristique $p$; leur produit $X\times X'$ est identique au produit fibré $X\times_eX'$, où $e=\Spec(\mathbb F_p)$. Il est clair que
\[
\fr_{X\times X'}=\fr_X\times\fr_{X'}=(\fr_X\times\id_{X'})\circ(\id_X\times\fr_{X'})
=(\id_X\times\fr_{X'})\circ(\fr_X\times\id_{X'}).
\]
On en déduit immédiatement que, pour tout $X$-préschéma $Y$ et tout $X'$-préschéma $Y'$, on a un isomorphisme canonique
\[
(Y\times Y')^{(p/X\times X')}\\simeq Y^{(p/X)}\times Y'{}^{(p/X')},
\]
modulo cet isomorphisme on peut écrire:
\[
\Fr_{Y\times Y'/X\times X'}=\Fr_{Y/X}\times\Fr_{Y'/X'}.
\]

\begin{statement}{Proposition 3}
Soient $X$ et $X'$ des préschémas de caractéristique $p$, $\Lambda$ un anneau commutatif, $F$ un faisceau de $\Lambda$-modules sur $X_{\et}$ et $F'$ un faisceau de $\Lambda$-modules sur $X'_{\et}$. On a
\[
\Fr^*_{F\boxt_\Lambda F'/X\times X'}
=\Fr^*_{F/X}\otimes\Fr^*_{F'/X'}
=(\Fr^*_{F/X}\otimes\id_{\fr_{X'}^*F'})\circ
(\id_{\fr_X^*F}\otimes\Fr^*_{F'/X'})
\]
\[
=(\id_F\otimes\Fr^*_{F'/X'})\circ
(\Fr^*_{F/X}\otimes\id_{F'}).
\]
\end{statement}

Pour écrire ces égalités, on a identifié
\[
\fr_{X\times X'}^*(F\boxt_\Lambda F')
\]
au produit tensoriel $\fr_X^*(F)\otimes_\Lambda\fr_{X'}^*(F')$ au moyen de l'isomorphisme canonique.

Les deux membres de la première égalité dépendent fonctoriellement de $F$ et de $F'$ et ils sont compatibles avec le changement de la base $X$ ou de la base $X'$; de plus ils sont compatibles avec le passage à la limite inductive sur $F$ ou sur $F'$, car le produit tensoriel commute à la limite inductive et il est clair que
\[
\Fr^*_{(\varinjlim F_i)/X}=\varinjlim\Fr^*_{F_i/X};
\]
ceci permet de se ramener au cas où $F=X$ et $F'=\Lambda_{X'}$, alors c'est évident.

\begin{statement}{Corollaire}
Avec les données de la proposition 3, les endomorphismes de $\RGamma_{X\times X'}(F\boxt F')$ définis respectivement par
\[
(\id_X\times\fr_{X'},\,\id_F\otimes\Fr^*_{F'/X'})
\]
et par
\[
(\fr_X\times\id_{X'},\,\Fr^*_{F/X}\otimes\id_{F'})
\]
sont des isomorphismes réciproques l'un de l'autre.
\end{statement}

En désignant ces endomorphismes par $\psi'$ et par $\psi$, montrons par exemple que $\psi\circ\psi'=\id$; cela résulte du diagramme commutatif suivant:
\[
\begin{tikzcd}[column sep=huge,row sep=large]
\RGamma_{X\times X'}(F\boxt F')
  \arrow[r,"\id_X\times\fr_{X'}"] \arrow[dr,"\fr_{X\times X'}"']
& \RGamma_{X\times X'}(F\boxt\fr_{X'}^*F')
  \arrow[r,"\RGamma(\id_F\otimes\Fr^*_{F'/X'})"] \arrow[d,"\fr_X\times\id_{X'}"]
& \RGamma_{X\times X'}(F\boxt F') \arrow[d,"\fr_X\times\id_{X'}"]\\
& \RGamma_{X\times X'}(\fr_X^*F\boxt\fr_{X'}^*F')
  \arrow[r,"\RGamma(\Fr^*_{F/X}\otimes\id_{F'})"']
& \RGamma_{X\times X'}(F\boxt F') ,
\end{tikzcd}
\]
qui est commutatif à cause de
\[
\fr_{X\times X'}=(\fr_X\times\id_{X'})\circ(\id_X\times\fr_{X'})
\]
pour le triangle du haut et de
\[
\Fr^*_{F\boxt F'/X\times X'}=(\Fr^*_{F/X}\otimes\id_{F'})\circ(\id_{\fr_X^*F}\otimes\Fr^*_{F'/X'})
\]
pour le triangle du bas, et dans lequel le composé des deux flèches obliques est l'identité d'après la proposition 2c). À l'aide d'un diagramme analogue échangeant les rôles de $(X,F)$ et $(X',F')$, on obtient $\psi'\circ\psi=\id$.

Nous appliquerons ce résultat au cas particulier où $X'=e=\Spec(\overline{\mathbb F}_p)$ est le spectre d'une clôture algébrique $\overline{\mathbb F}_p$ de $\mathbb F_p$ et où $F'=\Lambda_e$ est le faisceau constant de fibre $\Lambda$. Dans ce cas nous poserons
\[
\overline X=X\times_e e=X\otimes_{\mathbb F_p}\overline{\mathbb F}_p,
\qquad
\overline F=F\otimes_\Lambda\Lambda_e,
\]
image réciproque de $F$ par la projection $\overline X\to X$. L'endomorphisme de Frobenius $\fr_e$ s'identifie au générateur canonique
\[
f:\lambda\longmapsto\lambda^p\qquad(\lambda\in\overline{\mathbb F}_p)
\]
du groupe de Galois
\[
\pi_e=\Gal(\overline{\mathbb F}_p/\mathbb F_p),
\]
$f$ déterminant un isomorphisme $\widehat{\mathbb Z}\simeq\pi_e$, et
\[
\Fr^*_{\Lambda_e/e}=\id_{\Lambda_e}
\]
car $\Lambda_e$ est un faisceau constant. Posons
\[
\overline{\fr}_X=\fr_X\times\id_e,
\]
image réciproque de $\fr_X$ par $\overline X\to X$,
\[
f_X=\id_X\times\fr_e,
\]
automorphisme de $\overline X$ défini par $f$ par transport de structure, et
\[
\overline{\Fr}^*_{F/X}=\Fr^*_{F/X}\otimes\id_{\Lambda_e},
\]
image réciproque de $\Fr^*_{F/X}$ par $\overline X\to X$. On peut écrire:
\[
\fr_{\overline X}=\overline{\fr}_X\circ f=f\circ\overline{\fr}_X,
\qquad
\Fr^*_{\overline F/\overline X}=\overline{\Fr}^*_{F/X}.
\]
Le corollaire de la proposition 3 montre que l'endomorphisme de $\RGamma_{\overline X}(\overline F)$ défini par le couple $(\fr_X,\overline{\Fr}^*_{F/X})$ est l'automorphisme inverse de celui que définit $f_X$. Nous dirons que $(\fr_X,\overline{\Fr}^*_{F/X})$ est la correspondance de Frobenius géométrique, c'est l'image inverse par $\overline X\to X$ de la correspondance de Frobenius sur $(X,F)$, et nous désignerons par
\[
\fr_{\RGamma_{\overline X}(\overline F)}
\]
l'automorphisme de $\RGamma_{\overline X}(\overline F)$ qu'elle définit; l'opération de Frobenius arithmétique, provenant par transport de structure de $f\in\pi_e$, en est l'inverse.

\section*{§ 3. La fonction $L$}

\subsection*{1. Définition}

Soit $p$ un nombre premier; considérons un nombre premier $\ell\ne p$ et une extension finie $\Omega$ de $\mathbb Q_\ell$. À tout couple $(X,F)$ d'un schéma $X$ de type fini sur $e=\Spec(\mathbb F_p)$ et d'un $\Omega$-faisceau constructible $F$ sur $X$, nous allons associer une série formelle
\[
L_F\in\Omega[[t]],
\]
de terme constant $1$.

\end{document}
